\section{Compilazione e configurazione dell’ambiente}

La sequenza di build adottata è riportata in questa sezione in forma riproducibile. L’ordine è importante perché alcuni progetti richiedono artefatti compilati da dipendenze precedenti.

\subsection{OMNeT++ 6.2}

Per OMNeT++ è stato seguito il flusso ufficiale basato su \texttt{install.sh} e successiva compilazione locale \cite{omnet_download,omnet_versions}. Un esempio minimale è riportato di seguito.

\begin{lstlisting}[caption={Installazione e build OMNeT++ 6.2},label={lst:omnet_build}]
tar -xvf omnetpp-6.2.0-src-linux.tgz
cd omnetpp-6.2.0

# Installa dipendenze e crea il virtualenv Python
# Esegue in automatico anche configurazione, compilazione e importazione delle variabili d'ambiente
./install.sh
\end{lstlisting}

\subsection{SUMO 1.20 con dipendenze aggiornate}

Il workflow usato in tesi è stato:
\begin{enumerate}
\item installare le dipendenze indicate dal repository SUMO corrente;
\item fare checkout della versione 1.20;
\item compilare da sorgente con CMake.
\end{enumerate}

\begin{lstlisting}[caption={Build SUMO 1.20 con dipendenze dal repository ufficiale},label={lst:sumo_build}]
git clone https://github.com/eclipse-sumo/sumo.git
cd sumo

# Dipendenze (repository ufficiale SUMO)
sudo apt-get install \$(cat build_config/build_req_deb.txt build_config/tools_req_deb.txt)

# Passaggio alla versione 1.20 (v1.20 nome arbitrario per il branch creato)
git checkout -b v1.20 v1\_20\_0

# Compilazione
cmake -B build .
cmake --build build -j$(nproc)

\end{lstlisting}

\subsection{INET, Simu5G, Veins, Plexe e sotto-progetti}

Per gli altri componenti è stata seguita la documentazione ufficiale, con particolare riferimento alle istruzioni di \texttt{plexe\_5g} per la catena INET $\rightarrow$ Simu5G $\rightarrow$ Veins $\rightarrow$ Plexe \cite{plexe_5g_readme,plexe_building}.

\begin{lstlisting}[caption={Sequenza di compilazione dei framework OMNeT++ based},label={lst:framework_build}]
# INET
git clone https://github.com/inet-framework/inet
cd inet/
git checkout -b inet_v4.5.4 v4.5.4
source setenv

make makefiles
make -j <ncores> MODE=release
make -j <ncores> MODE=debug

# Simu5G (branch indicato da plexe_5g)
git clone https://github.com/michele-segata/Simu5G/ simu5g
cd simu5g/
git checkout plexe_mec
source setenv

make makefiles
make -j <ncores> MODE=release
make -j <ncores> MODE=debug

# Veins (+ eventuale sotto-progetto veins_inet)
git clone https://github.com/sommer/veins/
cd veins/
git checkout -b veins_v5.3.1 veins-5.3.1
source setenv

./configure
make -j <ncores>

# veins_inet
cd subprojects/veins_inet/
./configure
make -j <ncores>

# Plexe + subprojects
git clone https://github.com/michele-segata/plexe
cd plexe/
git checkout -b plexe_v3.2 plexe-3.2
source setenv

./configure
make -j <ncores>

cd subprojects/plexe\_5g
./configure
make -j <ncores>

\end{lstlisting}